% Pursuing Stacks (cotes 134-2 to 134-8): archive cote 134-6, site page 98 (PDF page 99), typescript page 416.
% typescript. Draft compared: ps5.tex:6532-6599.
% AI-transcribed by Claude Opus 5.5 (claude-opus-5-5, adaptive thinking, effort high), Message Batches API, batch ps3 (msgbatch_01XdFmuVVBAcgc7EFxdtCqsj), 2026-11-12; not yet checked by eye.
\scanpage{134-6/98}
\mspage{416}
deduce from this the expression
\[
(83)\qquad \mathrm{Hom}(L,F) = \overset{\hand{$\sim$}}{F}(L)\hand{$\simeq$}\ \underset{\hand{(x,u)}\ \text{in Add}(A)_{/L}}{\underset{\hand{$\longleftarrow$}}{\lim}}\ \xout{\unsure{Add(x)}}\ \mathrm{Add}(F)(x) \qquad ,
\]
which in an evident way is functorial in L for variable L.

This is about the best which can be done in general,  it seems to me,
by way of ``computational'' expression of Hom(L,F) in terms of F and L
given as in (79).  Of course the symbol Add(A)$_{/L}$ is relative to the cano-%
nical functor \hand{$\mathcal{E}$}$_{\mathrm{ab}}$ (80),  which is a full embedding in case the site
structure on A is the chaotic one i.e. \hand{$\mathcal{A}$}$=$ A\^{}.  In computational terms,
this category is rather explicit,  an object of the category is just a
pair
\[
\xout{\unsure{$((a_i)_{i\in I},(u_i)_{i\in I})$}}\ ( x , u) = ((x_i)_{i\hand{$\in$} I} , (u_i)_{i\hand{$\in$} I})
\]
where I is a finite indexing set,  $(a_i)_{i\hand{$\in$} I}$ a family of objects of A,  and
for i in I\hand{,} u$_i$ is an element of  L(a$_i$)  -  I'll leave to the reader the
description of maps between such objects.  The value of Add(F)(x) is given
by (81) above.

\uline{Remark}.  The expression (83) of Hom(L\xout{x},F) $=\overset{\hand{$\sim$}}{F}$(L)  \inserted[typed]{makes sense}\xout{\unsure{has a meaning}}, provided
only the additive category \xout{xx}\hand{$\mathcal{M}$} is stable under small inverse limits,
without having to assume that \hand{$\mathcal{M}$} be a pseudotopos.  This makes us suspect
that \xout{the fact that} \hand{the functor}
\[
\hand{$\unsure{\Phi}\mapsto\Phi$}\circ\mathrm{Wh} \ :\qquad \uline{Hom}^{\hand{!}}(\hand{$\mathcal{A}$}^{o}_{ab},\hand{$\mathcal{M}$})\hand{$\longrightarrow\mathcal{A}_{\mathcal{M}}$}
\]
is an equivalence ((13)  p.388)  \xout{is valid} without \xout{havin} this extra assump-%
tion,  provided \hand{$\mathcal{A}$} is an actual topos (not only a pseudotopos as in loc.
cit.)\hand{,} \inserted[typed]{Indeed} \deleted{\xout{xx}} we got a reasonable candidate for a quasi-inverse functor
\[
F\hand{$\mapsto$}\overset{\hand{$\sim$}}{F}\ :\quad \hand{$\mathcal{A}_{\mathcal{M}}\longrightarrow$}\ \uline{Hom}^{\hand{!}}(\hand{$\mathcal{A}$}_{ab}{}^{o},\hand{$\mathcal{M}$}) \qquad .
\]
The only point \inserted[typed]{still} to check,  with (83) defining $\overset{\hand{$\sim$}}{F}$ for given F in \hand{$\mathcal{A}_{\mathcal{M}}$},  is
that we get a functorial isomorphism
\[
\overset{\hand{$\sim$}}{F}\ (\ \underline{Z}^{(a)})\ \hand{$\simeq$}\ F(a)
\]
for a in A.  In case \hand{$\mathcal{A}$}$=$ A\^{} this follows from the fact that (8\over{3}{0}) is fully
faithful,  hence  Add(A)$_{/L}$  for  L $=\underline{Z}^{(a)}$ admits a final object  -  hence
the limit (83) is the value of Add(F) on the latter,  namely F(a).

I feel the little program on the \uline{Hom} and Hom  operations\hand{,} as contempla-%
ted on page 412\hand{,} is by now completed;  all we've got to do still is to
dualize \xout{the results} to get corresponding  results for $\ast$ and \hand{$\circledast$}.  It's
just a matter of essentially copying the formulaire (75),  which I'll \hand{do}

\margin{98}
